\addchap{Notation}

The chapters share the notation below. A chapter that needs a clashing local symbol says so where it introduces it, and the notes at the end of each chapter translate the notation of its main sources into this one.

\addsec*{Parties and protocols}

\begin{tabular}{@{}p{0.27\textwidth}p{0.69\textwidth}@{}}
$\V$ & the verifier; in a debate, the judge \\
$\Prov$ & the prover, and the honest prover's strategy \\
$\Prov^*$ & an arbitrary prover, which may cheat \\
$\Prov_1,\Prov_2$ & two provers questioned separately \\
$\party{A},\party{B}$ & two debaters \\
$\party{H}$ & a human judge consulted as an oracle \\
$x$, $n=\abs{x}$ & the input and its length \\
$L$ & a language: the set of inputs for which the claim is true \\
$\interact{\Prov}{\V}(x)$ & the verifier's output, \acc{} or \rej, after interacting with $\Prov$ on $x$ \\
$c$, $s$ & completeness and soundness error; $1-s$ is the probability that a lie is caught \\
$k$, $q$, $t$ & rounds, queries to a proof, repetitions \\
$\rho$ & the verifier's random bits; $r$ a random field element \\
$T$ & the running time of the computation that a claim is about \\
$\lambda$, $\negl(\lambda)$ & a security parameter and a negligible function of it
\end{tabular}

\addsec*{Classes}

Complexity classes are set in bold: $\cc{P}$, $\cc{NP}$, $\cc{coNP}$, $\cc{BPP}$, $\cc{MA}$, $\cc{AM}$, $\cc{IP}$, $\cc{PSPACE}$, $\cc{EXP}$, $\cc{NEXP}$, $\cc{MIP}$, $\cc{\#P}$, and $\cc{PCP}(r,q)$ for languages with probabilistically checkable proofs that use $r$ random bits and $q$ queries. The appendix defines most of them; \cref{ch:spotchecks} defines $\cc{PCP}(r,q)$, and \cref{ch:crossexam} defines $\cc{MIP}$.

\addsec*{Algebra and probability}

\begin{tabular}{@{}p{0.27\textwidth}p{0.69\textwidth}@{}}
$\F$, $\abs{\F}$ & a finite field and its size \\
$\ell$, $d$ & the number of variables of a polynomial and its degree in each variable \\
$\tilde f$ & the multilinear extension of a function $f\colon\bits^\ell\to\F$ \\
$S$, $D$ & the size and depth of a circuit \\
$\bits$ & the set $\{0,1\}$; $\bits^n$ the strings of length $n$ \\
$\chi_S$, $\hat F(S)$ & the character $x\mapsto(-1)^{\sum_{i\in S}x_i}$ of $\F_2^n$, and the Fourier coefficient of $F=(-1)^f$ at $S$ \\
$\Prob$, $\E$ & probability and expectation, always over stated randomness \\
$m$, $\cD$, $\cH$ & a number of samples, a distribution, a hypothesis class \\
$\eps$, $\delta$ & the accuracy and the confidence of a learning guarantee
\end{tabular}

\addsec*{Asymptotics}

For functions of the input length $n$ (or of $\lambda$), $f=O(g)$ means $f\le Cg$ for a constant $C$ and all large $n$; $\Omega$, $\Theta$ and $o$ are as usual, $\poly(n)$ is some polynomial in $n$, and $\polylog(n)$ some polynomial in $\log n$. Logarithms are to base $2$ unless written $\ln$.
